\begin{frame}{\insertsubsection}

  CVE-2018-3639: спекулятивное выполнение чтения памяти после сохранения её в
  регистр

  \begin{figure}[h]
    \includegraphics[height = .7\textheight]{spectre_logo}
    \caption{Spectre}
  \end{figure}

  \note{


  }
\end{frame}

\subsubsection{Сначала чтение, потом запись}
\begin{frame}[fragile]{\insertsubsubsection}
  \begin{minted}{nasm}
    STR X1, [X2]   ; X2 - адрес памяти, который ещё не известен
    ...
    LDR X3, [X4]   ; X4 содержит тот же адрес, что и X2
    <произвольная обработка X3>
    LDR X5, [X6, X3]  ; спекулятивное выполнение со старым адресом
  \end{minted}

  \note{

    Во многих современных процессорах применяются интересные техники
    оптимизации, а именно: \textbf{загрузка данных} из памяти по определённому
    адресу производится \textbf{раньше, чем запись} в тот же участок памяти в
    случае, \textbf{если адрес} на этапе записи данных \textbf{ещё не известен}.

    В итоге у нас \textbf{спекулятивно выполняются} все операции \textbf{со
      старым адресом} в регистре (в том числе запись в кэш).

  }
\end{frame}

\subsubsection{Читаем данные EL1}
\begin{frame}[fragile]{\insertsubsubsection}
  \begin{minted}{nasm}
    STR X1, [X2]
    ...
    ERET           ; возврат на более нижний уровень исключений
    ...
    LDR X3, [X4]   ; X4 содержит такой же физический адрес, как и X2,
                   ; но виртуальный адрес отличается
    <произвольная обработка X3>
    LDR X5, [X6, X3]
  \end{minted}

  \note{

    Если \textbf{адрес один и тот же} (виртуальный и физический), то данная
    уязвимость эксплуатируема только \textbf{на одном уровне исключений}
    (exception level for ARM).

    В случае, если есть возможность \textbf{чтения на одном уровне исключений},
    а \textbf{загрузки на другом}, и при этом используется \textbf{один и тот же
      физический адрес}, то есть возможность эксплуатации уязвимости \textbf{на
      разных уровнях}.

  }
\end{frame}

\subsubsection{Спекулятивное чтение одного и того же регистра}
\begin{frame}[fragile]{\insertsubsubsection}
  \begin{minted}{nasm}
    STR X1, [SP]
    ...
    LDR X3, [SP]
    <произвольная обработка X3>
    LDR X5, [X6, X3]
    <произвольная обработка X5>
    LDR X7, [X8, X5]
  \end{minted}

  \note{

    Удивительно, но факт --- мы можем читать данные из памяти, обращаясь по
    одному и тому же регистру.

    Такое возможно благодаря \textbf{современным оптимизациям на некоторых
      архитектурах}. В случае, если у нас содержимое SP, например,
    \textbf{отсутствовало в кэше}, то \textbf{при записи} мы получим
    \textbf{cache miss} и задержку, позволяющую нам \textbf{спекулятивно
      прочитать всё те же данные}. Такое возможно по причине того, что процессор
    в RoB отслеживает, когда данные из регистра попали в кэш и соответственно
    ускоряет исполнение инструкций.

    \textbf{PoC нет!}

  }
\end{frame}

\subsubsection{Спекулятивный запуск непривилегированного кода}
\begin{frame}[fragile]{\insertsubsubsection}
  \begin{minted}{nasm}
    STR X1, [SP]
    ...
    LDR X3, [SP]
    ...
    BLR X3
  \end{minted}

  \note{

    Так же, как и в обычном Spectre, мы можем составить цепочку ROP гаджетов и
    спекулятивно их запустить, чтобы \textbf{записать нужные нам данные в кэш}.

  }
\end{frame}

\subsubsection{Сначала запись, потом чтение}
\begin{frame}[fragile]{\insertsubsubsection}
  \begin{minted}{nasm}
    ...
    LDR X3, [X4]
    <произвольная обработка X3>
    LDR X5, [X6, X3]
    ....
    STR X1, [X2] ; X2 содержит тот же адрес, что и X4
  \end{minted}

  \note{

    Утверждается, что существует возможность прочитать данные спекулятивно,
    записанные также спекулятивно, но позже. Это возможно в случае, если регистр
    для чтения будет высчитываться гораздо дольше регистра для записи.

    \textbf{PoC нет!}

  }
\end{frame}

\subsubsection{Предотвращение}
\begin{frame}{\insertsubsubsection}

  \LARGE

  \begin{itemize}
  \item вставка «барьеров»
  \item отключение реорганизации операций чтения и записи
  \item SafeSpec
  \end{itemize}

  \note{

    Медленно!

  }

\end{frame}
